\chapter{CW Complexes}
\label{ch:viii06-cw-complexes}

Cell attachments become a general language when organized into skeleta.  CW
complexes are built from disks in increasing dimension, together with closure
finiteness and the weak topology.  They include graphs, spheres, surfaces,
projective spaces, and most spaces used in classical algebraic topology.

\section{Cells and characteristic maps}
An \(n\)-cell is homeomorphic to \(\mathring D^n\cong\mathbb R^n\).  A
characteristic map
\[
\chi_\alpha:D^n\to X
\]
restricts to a homeomorphism from \(\mathring D^n\) onto its open cell and sends
the boundary into lower-dimensional cells.

\section{Skeleta}
\begin{definition}[Skeleton]\label{def:viii06-skeleton}
The \(n\)-skeleton \(X^n\) is the union of all cells of dimension at most \(n\),
with \(X^{-1}=\varnothing\).
\end{definition}
Inductively, \(X^n\) is obtained from \(X^{n-1}\) by attaching \(n\)-cells.

\section{CW complex}
\begin{definition}[CW complex]\label{def:viii06-cw}
A CW complex is a cell decomposition satisfying:
\begin{enumerate}
\item every cell has a characteristic map with boundary in the lower skeleton;
\item closure finiteness: the closure of each cell meets only finitely many cells;
\item weak topology: \(A\subset X\) is closed exactly when
\(A\cap\overline e\) is closed in \(\overline e\) for every cell \(e\).
\end{enumerate}
\end{definition}
The letters C and W refer to closure finiteness and weak topology.

\section{Skeletal pushouts}
At stage \(n\),
\[
\begin{array}{ccc}
\coprod_\alpha S^{n-1}_\alpha&\longrightarrow&X^{n-1}\\
\big\downarrow&&\big\downarrow\\
\coprod_\alpha D^n_\alpha&\longrightarrow&X^n
\end{array}
\]
is a pushout.  This is VIII/05 repeated dimension by dimension.

\section{Dimension}
\begin{definition}[Dimension]\label{def:viii06-dimension}
\[
\dim X=\sup\{n:X\text{ has an \(n\)-cell}\}.
\]
\end{definition}
A zero-dimensional CW complex is discrete; a one-dimensional CW complex is a
topological graph.

\section{Finite and locally finite complexes}
\begin{definition}[Finite CW complex]\label{def:viii06-finite}
A CW complex is finite if it has only finitely many cells.
\end{definition}

\begin{proposition}[Finite implies compact]\label{prop:viii06-compact}
Every finite CW complex is compact.
\end{proposition}

\begin{definition}[Locally finite]\label{def:viii06-locally-finite}
A CW complex is locally finite if every point has a neighborhood meeting only
finitely many cells.
\end{definition}

\begin{warning}[Different finiteness notions]\label{warn:viii06-finiteness}
Closure finiteness and local finiteness are distinct conditions.
\end{warning}

\section{Subcomplexes}
\begin{definition}[Subcomplex]\label{def:viii06-subcomplex}
A subspace \(A\subset X\) is a CW subcomplex if it is a union of cells and,
with every cell it contains, it contains the closure of that cell.
\end{definition}

Every skeleton is a subcomplex.

\begin{proposition}[Subcomplex closedness]\label{prop:viii06-subcomplex-closed}
Every subcomplex of a CW complex is closed.
\end{proposition}

\section{Cellular maps}
\begin{definition}[Cellular map]\label{def:viii06-cellular}
A map \(f:X\to Y\) of CW complexes is cellular if
\[
f(X^n)\subset Y^n
\]
for every \(n\).
\end{definition}

\section{Standard examples}
\begin{example}[Sphere]\label{ex:viii06-sphere}
\(S^n\) has one \(0\)-cell and one \(n\)-cell.
\end{example}

\begin{example}[Graph]\label{ex:viii06-graph}
Graphs are one-dimensional CW complexes with vertices as zero-cells and edge
interiors as one-cells.
\end{example}

\begin{example}[Torus]\label{ex:viii06-torus}
The torus has one \(0\)-cell, two \(1\)-cells, and one \(2\)-cell.
\end{example}

\begin{example}[Projective space]\label{ex:viii06-projective}
\(\mathbb RP^n\) has one cell in each dimension \(0,\dots,n\).
\end{example}

\section{CW structure is extra data}
The same underlying space may carry different CW structures.  For example,
\(S^1\) can have one zero-cell and one one-cell, or two zero-cells and two
one-cells.

\begin{warning}[Raw cell counts]\label{warn:viii06-cell-counts}
Individual cell counts depend on the CW structure.  Invariant combinations,
such as Euler characteristic for finite CW complexes, require proof of
independence from the decomposition.
\end{warning}

\section{Infinite examples}
The ray \([0,\infty)\) has a locally finite CW structure with zero-cells at
nonnegative integers and one-cells \((k,k+1)\).  Also,
\[
\mathbb RP^\infty=\bigcup_{n\ge0}\mathbb RP^n
\]
has one cell in every nonnegative dimension.

\section{Weak topology and cellwise continuity}
\begin{proposition}[Cellwise criterion]\label{prop:viii06-cellwise}
A function \(f:X\to Y\) from a CW complex is continuous if each composite
\[
f\circ\chi_\alpha:D^{n_\alpha}\to Y
\]
is continuous.
\end{proposition}

\section{Compact subsets}
\begin{theorem}[Compact-cell finiteness]\label{thm:viii06-compact-cells}
Every compact subset of a CW complex meets only finitely many open cells.
\end{theorem}

Consequently a compact image lies in a finite subcomplex after adjoining the
finitely many cell closures that it meets.

\section{Connectedness and the one-skeleton}
\begin{proposition}[One-skeleton detects components]\label{prop:viii06-connected}
For a nonempty CW complex, path components are already detected by the
one-skeleton.  In particular, \(X\) is path connected exactly when \(X^1\) is
path connected.
\end{proposition}

\section{CW pairs}
\begin{definition}[CW pair]\label{def:viii06-pair}
A CW pair \((X,A)\) consists of a CW complex \(X\) and a subcomplex \(A\subset X\).
\end{definition}

\section{Quotients by subcomplexes}
\begin{proposition}[CW quotient]\label{prop:viii06-quotient}
If \(A\subset X\) is a subcomplex, then \(X/A\) has a natural CW structure:
\(A\) becomes a zero-cell and cells not in \(A\) survive with their dimensions.
\end{proposition}

\section{Euler characteristic preview}
For a finite CW complex with \(c_n\) \(n\)-cells, define
\[
\chi(X)=\sum_n(-1)^n c_n.
\]
For the standard torus structure,
\[
\chi(T^2)=1-2+1=0,
\]
while
\[
\chi(S^n)=1+(-1)^n.
\]
Later homology proves this alternating sum is independent of the finite CW
structure.

\section{Why CW complexes are effective}
CW complexes combine constructive flexibility with strong control: compact
data live in finite subcomplexes, maps can often be made cellular, subcomplex
inclusions support homotopy extension, and homology can be computed from the
skeletal filtration and attaching maps.  VIII/07 and VIII/08 exploit exactly
these features.


\section{Exercises}

\begin{exercise}\label{exr:viii06-01}
What do C and W stand for?
\end{exercise}
\begin{hint}
Recall the defining global conditions.

% VIII-PEDAGOGY-HINT-v1.1 exr:viii06-01 append
Separate the finiteness assertion about each closed cell from the topological test using characteristic maps of all closed disks.
\end{hint}
\begin{solution}
Closure finite and weak topology.
\end{solution}

\begin{exercise}\label{exr:viii06-02}
Define \(X^n\).
\end{exercise}
\begin{hint}
Collect cells up to dimension \(n\).

% VIII-PEDAGOGY-HINT-v1.1 exr:viii06-02 append
Include every dimension \(0\) through \(n\), not merely the cells of exactly dimension \(n\).
\end{hint}
\begin{solution}
The union of all cells of dimension at most \(n\).
\end{solution}

\begin{exercise}\label{exr:viii06-03}
What is \(X^{-1}\)?
\end{exercise}
\begin{hint}
It is the initial skeleton.

% VIII-PEDAGOGY-HINT-v1.1 exr:viii06-03 append
Identify the starting object before any zero-cells have been attached; no points have yet been introduced.
\end{hint}
\begin{solution}
\(\varnothing\).
\end{solution}

\begin{exercise}\label{exr:viii06-04}
How is \(X^n\) built from \(X^{n-1}\)?
\end{exercise}
\begin{hint}
Use VIII/05.

% VIII-PEDAGOGY-HINT-v1.1 exr:viii06-04 append
Write one attaching map \(S^{n-1}_\alpha\to X^{n-1}\) per new cell and form the pushout of the disjoint union of boundary inclusions.
\end{hint}
\begin{solution}
Attach the \(n\)-cells along maps \(S^{n-1}\to X^{n-1}\).
\end{solution}

\begin{exercise}\label{exr:viii06-05}
Define CW dimension.
\end{exercise}
\begin{hint}
Take the largest cell dimension or supremum.

% VIII-PEDAGOGY-HINT-v1.1 exr:viii06-05 append
An infinite number of cells need not imply unbounded dimension; compare an infinite graph with projective space of unbounded dimension.
\end{hint}
\begin{solution}
\(\dim X=\sup\{n:X\text{ has an \(n\)-cell}\}\).
\end{solution}

\begin{exercise}\label{exr:viii06-06}
Describe a zero-dimensional CW complex.
\end{exercise}
\begin{hint}
It has only points.

% VIII-PEDAGOGY-HINT-v1.1 exr:viii06-06 append
Apply the weak-topology test to an arbitrary subset: its inverse image in each characteristic zero-disk is automatically open and closed.
\end{hint}
\begin{solution}
It is discrete.
\end{solution}

\begin{exercise}\label{exr:viii06-07}
Describe a one-dimensional CW complex.
\end{exercise}
\begin{hint}
Use vertices and edges.

% VIII-PEDAGOGY-HINT-v1.1 exr:viii06-07 append
Record each one-cell's two endpoints and allow them to coincide; this captures loops as well as ordinary edges.
\end{hint}
\begin{solution}
It is a topological graph.
\end{solution}

\begin{exercise}\label{exr:viii06-08}
Define a finite CW complex.
\end{exercise}
\begin{hint}
Count all cells.

% VIII-PEDAGOGY-HINT-v1.1 exr:viii06-08 append
Distinguish finitely many cells in total from finitely many cells in each dimension and from finite dimension.
\end{hint}
\begin{solution}
It has finitely many cells.
\end{solution}

\begin{exercise}\label{exr:viii06-09}
Why is a finite CW complex compact?
\end{exercise}
\begin{hint}
Use characteristic disks.

% VIII-PEDAGOGY-HINT-v1.1 exr:viii06-09 append
Map the finite disjoint union of closed characteristic disks onto the complex; use compactness of that union and continuity of the quotient map.
\end{hint}
\begin{solution}
It is a continuous image of a finite union of compact closed disks.
\end{solution}

\begin{exercise}\label{exr:viii06-10}
Define a subcomplex.
\end{exercise}
\begin{hint}
Require closure under cell closures.

% VIII-PEDAGOGY-HINT-v1.1 exr:viii06-10 append
Including an open cell obliges you to include all boundary cells in its closure, rather than only selected interior points.
\end{hint}
\begin{solution}
A union of cells containing the closure of every cell it contains.
\end{solution}

\begin{exercise}\label{exr:viii06-11}
Are skeleta subcomplexes?
\end{exercise}
\begin{hint}
Apply the definition.

% VIII-PEDAGOGY-HINT-v1.1 exr:viii06-11 append
The attaching boundary of a cell of dimension at most \(n\) lies in a lower skeleton, so check the subcomplex condition cell by cell.
\end{hint}
\begin{solution}
Yes.
\end{solution}

\begin{exercise}\label{exr:viii06-12}
Define a cellular map.
\end{exercise}
\begin{hint}
Respect dimensions.

% VIII-PEDAGOGY-HINT-v1.1 exr:viii06-12 append
The requirement concerns every skeleton simultaneously; it does not require individual cells to map bijectively to cells.
\end{hint}
\begin{solution}
\(f(X^n)\subset Y^n\) for all \(n\).
\end{solution}

\begin{exercise}\label{exr:viii06-13}
Give minimal cell counts for \(S^n\).
\end{exercise}
\begin{hint}
Use VIII/05.

% VIII-PEDAGOGY-HINT-v1.1 exr:viii06-13 append
Use \(D^n/\partial D^n\) for \(n>0\); check separately that \(S^0\) has two zero-cells.
\end{hint}
\begin{solution}
One \(0\)-cell and one \(n\)-cell.
\end{solution}

\begin{exercise}\label{exr:viii06-14}
Give torus cell counts.
\end{exercise}
\begin{hint}
Use the polygon model.

% VIII-PEDAGOGY-HINT-v1.1 exr:viii06-14 append
In the square quotient, count vertex equivalence classes and edge equivalence classes before the single open face.
\end{hint}
\begin{solution}
One \(0\)-cell, two \(1\)-cells, one \(2\)-cell.
\end{solution}

\begin{exercise}\label{exr:viii06-15}
Give \(\mathbb RP^n\) cell counts.
\end{exercise}
\begin{hint}
Use projective filtration.

% VIII-PEDAGOGY-HINT-v1.1 exr:viii06-15 append
At each inclusion \(\mathbb RP^{k-1}\subset\mathbb RP^k\), identify the complement as \(\mathbb R^k\).
\end{hint}
\begin{solution}
One cell in every dimension \(0,\dots,n\).
\end{solution}

\begin{exercise}\label{exr:viii06-16}
Can one space have different CW structures?
\end{exercise}
\begin{hint}
Think of \(S^1\).

% VIII-PEDAGOGY-HINT-v1.1 exr:viii06-16 append
Subdivide a circle edge by adding a vertex; compare its old and new vertex and edge counts without changing the underlying space.
\end{hint}
\begin{solution}
Yes.
\end{solution}

\begin{exercise}\label{exr:viii06-17}
Define local finiteness.
\end{exercise}
\begin{hint}
Look at neighborhoods.

% VIII-PEDAGOGY-HINT-v1.1 exr:viii06-17 append
State the quantifiers in the order: for each point, there exists a neighborhood, meeting only finitely many cells.
\end{hint}
\begin{solution}
Every point has a neighborhood meeting finitely many cells.
\end{solution}

\begin{exercise}\label{exr:viii06-18}
Is closure finiteness the same as local finiteness?
\end{exercise}
\begin{hint}
They quantify different things.

% VIII-PEDAGOGY-HINT-v1.1 exr:viii06-18 append
Use an infinite wedge of circles: each individual closed edge meets finitely many cells, while every neighborhood of the wedge point meets infinitely many edges.
\end{hint}
\begin{solution}
No.
\end{solution}

\begin{exercise}\label{exr:viii06-19}
State the cellwise continuity test.
\end{exercise}
\begin{hint}
Use characteristic maps.

% VIII-PEDAGOGY-HINT-v1.1 exr:viii06-19 append
Test the composite on each closed characteristic disk, including its boundary; continuity on open cells alone is insufficient.
\end{hint}
\begin{solution}
Continuity follows if the composite with every characteristic map is continuous.
\end{solution}

\begin{exercise}\label{exr:viii06-20}
What is true of compact subsets and cells?
\end{exercise}
\begin{hint}
Use compact-cell finiteness.

% VIII-PEDAGOGY-HINT-v1.1 exr:viii06-20 append
After identifying finitely many cells met by the compact set, use closure finiteness to place them inside a finite subcomplex.
\end{hint}
\begin{solution}
A compact subset meets only finitely many open cells.
\end{solution}

\begin{exercise}\label{exr:viii06-21}
Define a CW pair.
\end{exercise}
\begin{hint}
Choose a subcomplex.

% VIII-PEDAGOGY-HINT-v1.1 exr:viii06-21 append
The subspace must be a union of whole cells together with their boundary cells, not an arbitrary subset of a CW complex.
\end{hint}
\begin{solution}
A pair \((X,A)\) with \(A\) a subcomplex of CW complex \(X\).
\end{solution}

\begin{exercise}\label{exr:viii06-22}
Describe \(X/A\) cellularly.
\end{exercise}
\begin{hint}
Collapse the subcomplex.

% VIII-PEDAGOGY-HINT-v1.1 exr:viii06-22 append
For nonempty \(A\), first record the new collapsed vertex, then inspect the characteristic maps of cells outside \(A\).
\end{hint}
\begin{solution}
\(A\) becomes a zero-cell; other cells survive with their dimensions.
\end{solution}

\begin{exercise}\label{exr:viii06-23}
Compute \(\chi(T^2)\).
\end{exercise}
\begin{hint}
Use \(1-2+1\).

% VIII-PEDAGOGY-HINT-v1.1 exr:viii06-23 append
Use the alternating cell count with dimensions attached to each number; this prevents adding both edge contributions with the wrong sign.
\end{hint}
\begin{solution}
\(0\).
\end{solution}

\begin{exercise}\label{exr:viii06-24}
Compute \(\chi(S^n)\).
\end{exercise}
\begin{hint}
Use one \(0\)-cell and one \(n\)-cell.

% VIII-PEDAGOGY-HINT-v1.1 exr:viii06-24 append
The top cell contributes according to the parity of \(n\); distinguish the even- and odd-dimensional cases, including \(S^0\).
\end{hint}
\begin{solution}
\(1+(-1)^n\).
\end{solution}


\section{Solved problem dossiers}

\begin{problem}[Skeletal pushout]\label{prob:viii06-01}
Write the construction of \(X^n\) from \(X^{n-1}\).
\end{problem}
\begin{solution}
Take the pushout of \(\coprod S^{n-1}_\alpha\to X^{n-1}\) and the boundary inclusions \(\coprod S^{n-1}_\alpha\hookrightarrow\coprod D^n_\alpha\).
\end{solution}

\begin{problem}[Two circle structures]\label{prob:viii06-02}
Compare the one-vertex/one-edge and two-vertex/two-edge structures on \(S^1\).
\end{problem}
\begin{solution}
They have the same underlying space and different raw cell counts.  Both have Euler alternating sum zero.
\end{solution}

\begin{problem}[Sphere skeletons]\label{prob:viii06-03}
Describe the minimal CW filtration of \(S^n\).
\end{problem}
\begin{solution}
There is one zero-cell, no cells in dimensions \(1,\dots,n-1\), and one \(n\)-cell attached constantly.
\end{solution}

\begin{problem}[Torus skeleta]\label{prob:viii06-04}
Describe the standard torus skeleta.
\end{problem}
\begin{solution}
\(X^0\) is a point, \(X^1=S^1\vee S^1\), and \(X^2\) attaches one disk along \(aba^{-1}b^{-1}\).
\end{solution}

\begin{problem}[Projective skeleta]\label{prob:viii06-05}
Identify the \(k\)-skeleton of standard \(\mathbb RP^n\).
\end{problem}
\begin{solution}
It is \(\mathbb RP^k\), because the projective filtration adds exactly one cell at each dimension.
\end{solution}

\begin{problem}[Finite implies compact]\label{prob:viii06-06}
Give the compactness proof for a finite CW complex.
\end{problem}
\begin{solution}
Take the finite disjoint union of all compact characteristic disks and map it continuously and surjectively to the complex.
\end{solution}

\begin{problem}[Subcomplex closure]\label{prob:viii06-07}
Why does a subcomplex contain cell boundaries?
\end{problem}
\begin{solution}
By definition it contains the closure of each cell it contains, so all cells appearing in those boundaries are included.
\end{solution}

\begin{problem}[Skeleton subcomplex]\label{prob:viii06-08}
Prove \(X^n\) is a subcomplex.
\end{problem}
\begin{solution}
The closure of a \(k\)-cell with \(k\le n\) lies in cells of dimension at most \(k\), hence remains in \(X^n\).
\end{solution}

\begin{problem}[Cellular map test]\label{prob:viii06-09}
For minimal sphere CW structures, when is \(f:S^n\to S^n\) cellular?
\end{problem}
\begin{solution}
It suffices that the unique zero-cell maps to the unique zero-cell; the full \(n\)-skeleton condition is automatic.
\end{solution}

\begin{problem}[Graph realization]\label{prob:viii06-10}
Explain how a multigraph becomes a CW complex.
\end{problem}
\begin{solution}
Use vertices as zero-cells and attach one interval per edge to its endpoint vertices, allowing loops and parallel edges.
\end{solution}

\begin{problem}[Infinite ray]\label{prob:viii06-11}
Verify the integer subdivision of \([0,\infty)\) is locally finite.
\end{problem}
\begin{solution}
A small neighborhood of a nonvertex meets one edge; near an integer it meets at most that vertex and its adjacent edges.
\end{solution}

\begin{problem}[Finiteness distinction]\label{prob:viii06-12}
Contrast closure finiteness and local finiteness.
\end{problem}
\begin{solution}
Closure finiteness starts from one cell and studies its closure; local finiteness starts from one point and studies a neighborhood.
\end{solution}

\begin{problem}[Weak topology]\label{prob:viii06-13}
Why do compatible cellwise continuous maps assemble globally?
\end{problem}
\begin{solution}
Preimages of closed sets intersect every closed cell in closed sets; the weak topology then makes the full preimage closed.
\end{solution}

\begin{problem}[Compact image]\label{prob:viii06-14}
What does a continuous map from compact \(K\) into a CW complex reduce to?
\end{problem}
\begin{solution}
Its image meets finitely many cells and is contained in a finite subcomplex after adjoining their closures.
\end{solution}

\begin{problem}[CW quotient]\label{prob:viii06-15}
Describe cells of \(X/A\) for subcomplex \(A\).
\end{problem}
\begin{solution}
Collapse \(A\) to one zero-cell and retain each open cell outside \(A\), with the attaching map followed by the quotient.
\end{solution}

\begin{problem}[CW pair]\label{prob:viii06-16}
Give a natural example of a CW pair.
\end{problem}
\begin{solution}
With compatible cell structures, \((D^n,S^{n-1})\) is a CW pair because the boundary is a subcomplex.
\end{solution}

\begin{problem}[Euler sphere]\label{prob:viii06-17}
Compute Euler characteristics of even and odd spheres.
\end{problem}
\begin{solution}
\(\chi(S^{2k})=2\) and \(\chi(S^{2k+1})=0\).
\end{solution}

\begin{problem}[Euler projective]\label{prob:viii06-18}
Compute the cellular Euler sum for \(\mathbb RP^n\).
\end{problem}
\begin{solution}
\(\sum_{k=0}^n(-1)^k\), equal to \(1\) for even \(n\) and \(0\) for odd \(n\).
\end{solution}

\begin{problem}[One-skeleton connectedness]\label{prob:viii06-19}
Why can higher-dimensional cells not join two distinct components of the lower skeleton?
\end{problem}
\begin{solution}
For \(n\ge2\), the connected boundary sphere of an attached \(n\)-cell maps into a single component, so the disk is added to that component only.
\end{solution}

\begin{problem}[Bridge to mapping cones]\label{prob:viii06-20}
Why are CW attachments naturally related to mapping cones?
\end{problem}
\begin{solution}
Each disk is \(CS^{n-1}\); attaching it to \(X^{n-1}\) is the mapping cone of its sphere attaching map, the subject of VIII/07.
\end{solution}
